%==============================================================================
% Chapitre 10 - Densité de l'air
%==============================================================================

\section{Introduction}

La densité de l'air est l'un des paramètres les plus importants en
balistique extérieure.

Elle influence directement :

\begin{itemize}
\item la traînée ;
\item la conservation de la vitesse ;
\item le temps de vol ;
\item la portée ;
\item la dérive due au vent.
\end{itemize}

La densité de l'air dépend principalement :

\begin{itemize}
\item de la pression ;
\item de la température ;
\item de l'humidité.
\end{itemize}

De nombreux calculateurs balistiques utilisent directement cette grandeur.

\begin{aretenir}

La densité de l'air constitue l'un des paramètres atmosphériques les plus
importants pour le calcul des trajectoires.

\end{aretenir}

\section{Qu'est-ce que la densité ?}

La densité volumique, ou masse volumique, représente la quantité de matière
contenue dans un volume donné.

Elle est généralement notée :

\begin{center}$\rho$\end{center}

et s'exprime en :

\begin{center}$\mathrm{kg/m^3}$\end{center}

Plus la densité est élevée, plus la quantité d'air rencontrée par le
projectile est importante.

\section{Densité de référence}

Dans l'atmosphère standard ISA, au niveau moyen de la mer :

\begin{equation}
\rho_0 = 1,225~\mathrm{kg/m^3}
\end{equation}

Cette valeur constitue une référence fréquemment utilisée dans les calculs
balistiques.

\section{Influence de la température}

Lorsque la température augmente :

\begin{itemize}
\item l'air se dilate ;
\item son volume augmente ;
\item sa densité diminue.
\end{itemize}

À pression identique :

\begin{itemize}
\item l'air chaud est moins dense ;
\item l'air froid est plus dense.
\end{itemize}

\begin{exemple}

Un tir réalisé à 35°C rencontrera généralement moins de résistance
aérodynamique qu'un tir réalisé à 0°C.

\end{exemple}

\section{Influence de la pression}

Lorsque la pression augmente :

\begin{itemize}
\item davantage de molécules occupent le même volume ;
\item la densité augmente.
\end{itemize}

À température constante :

\begin{itemize}
\item une pression élevée augmente la densité ;
\item une pression faible réduit la densité.
\end{itemize}

Cette relation explique en partie les différences observées entre les tirs
réalisés à basse et haute altitude.

\section{Influence de l'humidité}

Comme vu au chapitre précédent, la vapeur d'eau possède une masse molaire
inférieure à celle de l'air sec.

Lorsque l'humidité augmente :

\begin{itemize}
\item une partie de l'air sec est remplacée ;
\item la densité diminue légèrement.
\end{itemize}

L'effet reste généralement plus faible que celui de la température ou de la
pression.

\section{Loi des gaz parfaits}

Les gaz atmosphériques peuvent être approximés par le modèle des gaz
parfaits.

\begin{rigueur}

La relation fondamentale est :

\begin{equation}
P = \rho R T
\label{eq:gaz_parfait}
\end{equation}

où :

\begin{itemize}
\item $P$ est la pression ;
\item $\rho$ est la densité ;
\item $R$ est la constante spécifique du gaz ;
\item $T$ est la température absolue.
\end{itemize}

\end{rigueur}

Cette relation montre que pression, température et densité sont
directement liées.

\section{Pourquoi la densité est-elle si importante ?}

Pour un projectile, la densité représente directement la quantité d'air à
traverser.

Une augmentation de la densité entraîne généralement :

\begin{itemize}
\item une augmentation de la traînée ;
\item une diminution de la vitesse résiduelle ;
\item une augmentation du temps de vol ;
\item une augmentation de la chute.
\end{itemize}

À l'inverse, une diminution de la densité réduit ces effets.

\section{Altitude-densité}

La notion d'altitude-densité est très utilisée en aéronautique et en
balistique.

Elle permet de résumer en une seule grandeur l'effet combiné :

\begin{itemize}
\item de l'altitude ;
\item de la température ;
\item de la pression.
\end{itemize}

Deux sites situés à la même altitude géométrique peuvent présenter des
altitudes-densité très différentes.

Pour le projectile, c'est souvent cette grandeur qui est la plus
représentative des conditions de vol.

\begin{exemple}

Une journée chaude en montagne peut produire une densité de l'air
comparable à celle rencontrée à une altitude beaucoup plus élevée.

\end{exemple}

\section{Influence sur les essais}

La densité de l'air est rarement mesurée directement.

Elle est généralement calculée à partir :

\begin{itemize}
\item de la température ;
\item de la pression ;
\item de l'humidité.
\end{itemize}

Cette valeur peut ensuite être utilisée pour :

\begin{itemize}
\item corriger les résultats ;
\item alimenter un calculateur balistique ;
\item comparer plusieurs campagnes d'essais.
\end{itemize}

\begin{simulation}

La plupart des moteurs balistiques utilisent la densité de l'air pour
calculer les forces aérodynamiques.

Une variation de quelques pourcents de la densité peut entraîner des
écarts mesurables à longue distance.

\end{simulation}

\section{Vers la traînée aérodynamique}

La densité de l'air n'agit pas directement sur le projectile.

Elle agit par l'intermédiaire des forces aérodynamiques.

La plus importante de ces forces est la traînée.

Le chapitre consacré à la traînée montrera comment la densité intervient
dans le calcul des efforts aérodynamiques.

\section{À retenir}

\begin{aretenir}

\begin{itemize}
\item La densité représente la masse d'air contenue dans un volume
donné.
\item Elle dépend principalement de la température, de la pression et
de l'humidité.
\item Une densité élevée augmente la traînée.
\item Une densité faible favorise la conservation de la vitesse.
\item La densité est un paramètre central des modèles balistiques.
\end{itemize}

\end{aretenir}

